%% kapa: exam=1 round=34 year=2023 period=1 form=A subject=economics subject_ko=경제학원론 seq=17 number=57
%% source: Q-Net 감정평가사 기출문제 (cst00309) artlSeq=5211789 file=raw/1008/5211789/2023년도 제34회 감정평가사 1차시험 1교시.pdf
%% conversion: kordoc 4.24.51 -> tools/gen1.py (auto); review=proofread-sonnet; flags=none
\begin{question}{57}
가격차별에 관한 설명으로 옳지 않은 것은?
\begin{choices}
  \item 극장에서의 조조할인 요금제는 가격차별의 한 예이다.
  \item 이부가격제(two-part pricing)는 가격차별 전략 중 하나이다.
  \item 제3급 가격차별을 가능하게 하는 조건 중 하나는 전매가 불가능해야 한다는 것이다.
  \item 제3급 가격차별의 경우 수요의 가격 탄력성이 상대적으로 작은 시장에서 더 낮은 가격이 설정된다.
  \item 제1급 가격차별에서는 소비자잉여가 발생하지 않는다.
\end{choices}
\end{question}
\answer{4}
